
Die Kommunikation und Zusammenarbeit im Projektteam Coasterbot war stets konstruktiv, entspannt, harmonisch, zuverlässig und fair. 

Aufgrund persönlicher Neigung konnten die technischen Aufgabenbereiche hervorragend in einen mechanischen (Dennis Borutta), einen elektronischen (Gereon Such) sowie einen algorithmischen Anteil inklusive Simulation (Stefan Etringer) aufgeteilt werden. In diesen Bereichen konnte jeder mit seinem persönlichen Vorwissen dem Team viel Einarbeitungszeit ersparen. 

Eine Herausforderung war das dislozierte Arbeiten an einer nicht normierten explorativen Hardwareplattform. Da nicht jedes Teammitglied alle notwendigen Werkzeuge zur Montage hatte wurde dies durchmischt und postalisch verschickt. Notwendige nachträgliche Änderungen an der Hardware waren so natürlich schwer nachzupflegen, anders als ein Softwareupdate, welches wir durch ein Github-Pull wieder hätten teilen könnten. Dies ist aber nicht der Kommunikation im Team, sondern der Lernkurve mit den Komponenten geschuldet, wie vorgehend beschrieben. 

An dieser Stelle möchte ich auch Herrn Dennis Borutta für seine initiale und konstante Initiative danken. Projektidee und Anfangsmotivation stammen aus seiner Feder und er hat es über die gesamte Projektlaufzeit geschafft, das Team fokussiert zusammenzuhalten. Vielen Dank dafür!

